% ============================================================
% Informe final de trabajo de grado
% Programa de Ingeniería Física, Universidad EAFIT, 2026-2
%
% Contenido mínimo de las directrices del programa (2026-2):
%   portada; tablas de contenido, figuras y tablas; resumen;
%   introducción con los objetivos aprobados; capítulos de
%   desarrollo; conclusiones y trabajos futuros; bibliografía.
% ============================================================
\documentclass[12pt,letterpaper,oneside]{report}

% ============================================================
% Preámbulo del informe final. Parte del .tex del anteproyecto
% (fuente, idioma, paquetes) y le cambia la página a ICONTEC.
% ============================================================

% ---------- Codificación e idioma ----------
\usepackage[utf8]{inputenc}
\usepackage[T1]{fontenc}
% es-tabla: "Tabla" e "Índice de tablas" en lugar de "Cuadro", que es
% como las nombran las directrices ("tablas de contenido, de figuras y
% de tablas").
\usepackage[spanish,es-nodecimaldot,es-tabla]{babel}
\usepackage{csquotes}

% ---------- Fuente tipo Calibri (la del anteproyecto) ----------
% ICONTEC solo *sugiere* Arial 12; se conserva la del anteproyecto.
\usepackage{carlito}
\renewcommand{\familydefault}{\sfdefault}

% ---------- Página según ICONTEC (NTC 1486, 6.a actualización, 2008) ----------
% Los cuatro márgenes a 3 cm, la variante simétrica que ICONTEC admite
% (la pide para impresión a doble cara), para que el texto quede
% centrado en la hoja. La variante para encuadernar es left=4cm,
% right=2cm. Número de página centrado a 2 cm del borde inferior
% (footskip = 3 cm - 2 cm); título de capítulo a 4 cm (ver titlesec).
\usepackage[letterpaper, top=3cm, left=3cm, right=3cm, bottom=3cm,
            footskip=1cm]{geometry}

% ---------- Interlineado ----------
% ICONTEC pide interlineado sencillo con una línea en blanco entre
% párrafos, que es una convención de Word. Se deja el 1.5 del
% anteproyecto, que es más legible con ecuaciones.
\usepackage{setspace}
\onehalfspacing

% ---------- Matemáticas, tablas y gráficos ----------
\usepackage{amsmath,amssymb}
\usepackage{array}
\usepackage{booktabs}
\usepackage{multirow}
\usepackage{graphicx}
\graphicspath{{figuras/}}
\usepackage{xcolor}
\usepackage{float}
\usepackage[font=small,labelfont=bf]{caption}
\usepackage{subcaption}

% ---------- Unidades ----------
\usepackage{siunitx}
\sisetup{output-decimal-marker={.}, group-digits=false}

% ---------- Listas compactas ----------
\usepackage{enumitem}
\setlist{itemsep=2pt, topsep=4pt, parsep=0pt}

% ---------- Títulos según ICONTEC ----------
% Capítulo en página nueva, centrado, en mayúscula sostenida y a 4 cm
% del borde (3 cm de margen + 1 cm).
\usepackage{titlesec}
\titleformat{\chapter}[block]
  {\normalfont\large\bfseries\centering}{\thechapter.}{0.5em}{\MakeUppercase}
\titlespacing*{\chapter}{0pt}{1cm}{1cm}
\titleformat{\section}{\normalfont\bfseries}{\thesection}{0.8em}{}
\titleformat{\subsection}{\normalfont\bfseries\itshape}{\thesubsection}{0.8em}{}
\titlespacing*{\section}{0pt}{12pt}{6pt}
\titlespacing*{\subsection}{0pt}{10pt}{4pt}

% ---------- Nombres de las listas ----------
\addto\captionsspanish{%
  \renewcommand{\contentsname}{Tabla de contenido}%
  \renewcommand{\listfigurename}{Lista de figuras}%
  \renewcommand{\listtablename}{Lista de tablas}%
}

% ---------- Bibliografía: IEEE numérico, como el anteproyecto ----------
% Para pasar a ICONTEC (NTC 5613) basta cambiar style=ieee.
% En Overleaf compila con biber sin configurar nada.
\usepackage[backend=biber, style=ieee]{biblatex}
\addbibresource{referencias.bib}

% ---------- Enlaces (va al final) ----------
\usepackage[hidelinks]{hyperref}

% ---------- Marcas de trabajo ----------
% \pendiente{...} deja en rojo lo que falta escribir. Antes de entregar,
% buscar "pendiente" en el proyecto debe dar cero resultados.
\newcommand{\pendiente}[1]{\textcolor{red}{[Pendiente: #1]}}


% ---------- Datos del trabajo ----------
\newcommand{\titulotrabajo}{Implementación y validación de propagadores
numéricos con muestreo controlado para la reconstrucción de hologramas en
microscopía holográfica digital sin lentes}
\newcommand{\autor}{Camilo Rodriguez Londoño}
\newcommand{\tutor}{Carlos Alejandro Trujillo Anaya}
\newcommand{\titulotutor}{PhD. Ciencias Físicas}
\newcommand{\fechaentrega}{2026} % TODO: mes de la entrega final (semana 15)

\hypersetup{pdftitle={\titulotrabajo}, pdfauthor={\autor}}

\begin{document}

% ---------- Preliminares ----------
% ICONTEC: la portada y los preliminares se cuentan pero no se numeran;
% el número se ve desde la introducción.
\pagenumbering{arabic}
% ============================================================
% PORTADA: la del anteproyecto, cambiando el tipo de trabajo.
% ICONTEC pide título, autor, tipo de trabajo y nombre y título
% académico del director, centrados y equidistantes.
% ============================================================
\begin{titlepage}
\centering
\vspace*{1cm}

{\large\bfseries UNIVERSIDAD EAFIT\par}
\vspace{0.3cm}
{\large ESCUELA DE CIENCIAS APLICADAS E INGENIERÍA\par}
\vspace{0.3cm}
{\large PROGRAMA DE INGENIERÍA FÍSICA\par}

\vspace{3cm}

{\Large\bfseries \titulotrabajo\par}

\vspace{2cm}

{\large Informe final de trabajo de grado\par}

\vspace{3cm}

{\large\textbf{Estudiante:}\par}
{\large \autor\par}
\vspace{0.5cm}
{\large\textbf{Tutor:}\par}
{\large \tutor\par}
{\normalsize \titulotutor\par}

\vfill

{\large Medellín, Colombia\par}
{\large \fechaentrega\par}
\end{titlepage}


\pagestyle{empty}
\assignpagestyle{\chapter}{empty} % también la primera hoja de cada lista
% ============================================================
% RESUMEN: se escribe al final, cuando los resultados estén.
% ============================================================
\chapter*{Resumen}

\pendiente{un párrafo (unas 250 palabras) con el problema, el método, los
resultados principales en números y la conclusión. Se escribe de último.}

\vspace{0.5cm}
\noindent\textbf{Palabras clave:} microscopía holográfica digital sin lentes;
espectro angular; producto matricial; muestreo; propagación numérica;
reconstrucción de hologramas.

% Abstract opcional: las directrices solo piden el resumen.
% \chapter*{Abstract}
% ...
% \noindent\textbf{Keywords:} digital lensless holographic microscopy;
% angular spectrum; matrix product; sampling; numerical propagation.


\tableofcontents
\listoffigures
\listoftables

% ---------- Cuerpo ----------
\pagestyle{plain}
\assignpagestyle{\chapter}{plain}
% ============================================================
% 1. INTRODUCCIÓN
% Base: secciones A, C, D y F del anteproyecto. Las directrices exigen
% los objetivos APROBADOS: van literales, no se reescriben.
% TODO: pasar a pasado lo que el anteproyecto dejaba en futuro.
% ============================================================
\chapter{Introducción}
\label{cap:introduccion}

% ------------------------------------------------------------
\section{Planteamiento del problema}
\label{sec:planteamiento}

En holografía digital el sensor no registra la imagen del objeto sino un patrón
de interferencia: para recuperarla hay que modelar numéricamente cómo viajó la
luz hasta el plano del objeto~\cite{buitrago2022}. Dentro de la teoría escalar
de la difracción ese cálculo ya tiene solución cerrada: la integral de
Rayleigh--Sommerfeld y el espectro angular (ASM) son la misma solución escrita
en el dominio espacial y en el de las frecuencias espaciales, y ninguna aproxima
el ángulo de propagación \cite{goodman2005, ersoy2006}. Lo que las separa es el
costo: el ASM se calcula con la transformada rápida de Fourier (FFT) en
$\mathcal{O}(N^2\log N)$ operaciones frente a las $\mathcal{O}(N^4)$ de la suma
directa, y por eso es el método más usado~\cite{buitrago2020}. La contrapartida
es que la exactitud deja de depender de la formulación y pasa a depender del
muestreo con que se discretiza.

La FFT impone sus propias reglas de muestreo: el espaciado de los puntos en
frecuencia queda fijado por el tamaño de la ventana de entrada, y el plano de
salida hereda el tamaño del de partida \cite{goodman2005}. Cuando esas reglas no
alcanzan a describir la propagación aparece \emph{aliasing}: frecuencias altas
que se confunden con bajas y contaminan la imagen. Las soluciones conocidas
pagan un precio: el método de banda limitada (BLAS) descarta las frecuencias
problemáticas y pierde detalle a distancias largas \cite{matsushima2009};
ampliar la ventana de entrada corrige el muestreo pero dispara memoria y tiempo
de cómputo \cite{yu2012}. Una alternativa reciente reemplaza la FFT por un
producto de matrices ---el método de espectro angular por producto matricial
(MPASM)~\cite{zhao2020}---: el muestreo deja de estar impuesto por el algoritmo
y se elige mediante dos coeficientes de control. $K_f$ comprime el intervalo de
muestreo en frecuencia, concentrando los puntos donde la función de
transferencia sí cumple la condición de muestreo; $s$ multiplica el número de
puntos en frecuencia y amplía el ancho de banda efectivo a distancias largas, a
costa de más operaciones.

La microscopía holográfica digital sin lentes (DLHM) es un escenario donde ese
control resulta pertinente. Al eliminar el objetivo del microscopio, la muestra
se ilumina con una onda esférica proveniente de un orificio muy pequeño
(\emph{pinhole}) y el sensor registra directamente el patrón de difracción, de
modo que la resolución queda determinada por la geometría del
montaje~\cite{lopera2024}. Dos rasgos de esa geometría tensionan al ASM
convencional: el paso del sensor deja el contenido frecuencial del campo cerca
del límite de Nyquist, y la iluminación divergente introduce una magnificación
que sitúa el patrón registrado en una escala distinta a la del objeto, algo que
la FFT no puede representar. El MPASM ofrece justamente ese control, pero solo
había sido evaluado en simulaciones con haces gaussianos e iluminación de onda
plana: no se había reportado su aplicación a hologramas experimentales ni
existían criterios publicados para elegir sus parámetros en condiciones reales
de registro.

\vspace{0.2cm}
\noindent\textbf{Pregunta de investigación:} ¿Qué características de un
propagador con control explícito del muestreo, como el MPASM, permiten mejorar
la calidad de las reconstrucciones en microscopía holográfica digital sin lentes
respecto de los propagadores empleados habitualmente en la técnica, y bajo qué
condiciones de registro resulta ventajoso?

% ------------------------------------------------------------
\section{Objetivos}
\label{sec:objetivos}

% Literales del anteproyecto aprobado. No editar.
\subsection*{Objetivo general}

Aplicar propagadores de campo óptico basados en el producto matricial, con
control flexible del muestreo espacial y frecuencial, a la reconstrucción de
hologramas en microscopía holográfica digital sin lentes.

\subsection*{Objetivos específicos}

\begin{enumerate}[label=\arabic*.]
  \item Implementar en un marco computacional unificado los métodos de
  propagación FFT-ASM, BLAS y MPASM, verificando su correcta ejecución mediante
  casos analíticos de referencia.

  \item Validar el desempeño de los métodos implementados bajo las condiciones
  de muestreo propias de DLHM, determinando los rangos de distancia, apertura
  numérica y coeficientes de control ($s$ y $K_f$, definidos en la
  Sección~\ref{sec:planteamiento}) en los que cada método mantiene una
  precisión aceptable, y contrastando la formulación matricial con la de la
  transformada Z \emph{chirp} (CZT), que permite un control de muestreo
  análogo.

  \item Comparar las reconstrucciones obtenidas con los propagadores
  implementados frente al método de espectro angular con iluminación de onda
  esférica, empleado actualmente en el montaje DLHM del laboratorio, mediante
  métricas cuantitativas de calidad sobre hologramas experimentales de muestras
  de resolución conocida.
\end{enumerate}

% ------------------------------------------------------------
\section{Justificación}

El MPASM aporta un grado de libertad que los métodos convencionales no tienen:
elegir de forma independiente cómo se muestrea el campo en el espacio y en la
frecuencia. Ese control encaja con lo que exige DLHM, donde la magnificación
geométrica y las condiciones de registro imponen restricciones que el ASM
basado en FFT no puede acomodar. Verificar si esa flexibilidad se traduce en
mejoras medibles sobre la reconstrucción de hologramas experimentales es una
pregunta que la literatura aún no respondía, ya que el método solo había sido
evaluado en simulaciones; y contar con criterios explícitos para elegir los
parámetros de muestreo evita el ajuste por ensayo y error que exige usar estos
propagadores. A esto se suma que su formulación matricial se parece
estructuralmente a la transformada Z \emph{chirp} (CZT), un algoritmo que
también permite calcular espectros con muestreo arbitrario; observar esa
relación durante la caracterización numérica ayuda a situar el método entre los
propagadores existentes.

% ------------------------------------------------------------
\section{Alcance}

El trabajo se limita al régimen de difracción escalar y a hologramas en línea
registrados en un montaje DLHM existente; no incluye instrumentación nueva, el
régimen vectorial ni algoritmos de recuperación de fase.

\pendiente{productos entregados: el paquete de software (nombre, licencia,
repositorio), los hologramas con sus reconstrucciones y los mapas de validez.
Decirlo en pasado y con lo que de verdad se entregó.}

% ------------------------------------------------------------
\section{Estructura del documento}

\pendiente{un párrafo que recorra los capítulos: el
Capítulo~\ref{cap:marco} presenta el marco teórico y los antecedentes; el
Capítulo~\ref{cap:metodos}, la implementación, los barridos y el montaje; el
Capítulo~\ref{cap:resultados}, los resultados por objetivo; y el
Capítulo~\ref{cap:conclusiones}, las conclusiones y el trabajo futuro.}

% ============================================================
% 2. MARCO TEÓRICO Y ESTADO DEL ARTE
% Lo que el lector necesita para entender los métodos y los
% resultados, y dónde está la literatura. Solo teoría ajena:
% lo que se desarrolló en el trabajo (p. ej. la distancia
% efectiva de la tarea 24) va en el capítulo 3.
% ============================================================
\chapter{Marco teórico y estado del arte}
\label{cap:marco}

% ------------------------------------------------------------
\section{Difracción escalar}
\label{sec:difraccion}

\pendiente{integral de Rayleigh--Sommerfeld y espectro angular como la misma
solución en dos dominios \cite{goodman2005, ersoy2006}; función de
transferencia $H(f_x,f_y)$; componentes evanescentes.}

% ------------------------------------------------------------
\section{Espectro angular por FFT y sus límites de muestreo}
\label{sec:asm-fft}

\pendiente{discretización; paso en frecuencia $1/L$ fijado por la ventana;
condición de muestreo de la función de transferencia y la frecuencia máxima
que la cumple (Ec.~12 de \cite{zhao2020}); aliasing; la ventana de salida
hereda la de entrada \cite{li2007}.}

% ------------------------------------------------------------
\section{Soluciones conocidas al submuestreo}
\label{sec:soluciones}

\subsection{Espectro angular de banda limitada (BLAS)}
\pendiente{recorte de la función de transferencia \cite{matsushima2009}; por
qué pierde detalle a distancias largas.}

\subsection{Ventana ampliada y muestreo no uniforme}
\pendiente{WWASM \cite{yu2012} y NUFFT \cite{kim2014}, en un párrafo cada uno:
qué resuelven y qué cuestan.}

% ------------------------------------------------------------
\section{Espectro angular por producto matricial (MPASM)}
\label{sec:mpasm}

\pendiente{la DFT como triple producto de matrices; los coeficientes $K_f$ y
$s$ y el muestreo de salida $r$; ancho de banda efectivo frente a BLAS;
complejidad \cite{zhao2020, zhao2020code}.}

\subsection{Relación con la transformada Z \emph{chirp}}
\pendiente{la CZT como otro cálculo de la DFT con muestreo arbitrario y el
parecido estructural con el MPASM (objetivo 2).}

% ------------------------------------------------------------
\section{Microscopía holográfica digital sin lentes}
\label{sec:dlhm}

\pendiente{geometría fuente puntual--muestra--sensor ($z$, $L$); magnificación
$M = L/z$; apertura numérica y resolución; modelo de formación del holograma
como difracción de un frente esférico más magnificación no homogénea
\cite{lopera2024}.}

\subsection{Reconstrucción en DLHM}
\pendiente{RS no aproximada como referencia de exactitud \cite{buitrago2019};
el espectro angular con iluminación de onda esférica que se usa de rutina en
el laboratorio, que es la línea base del objetivo 3.}

% ------------------------------------------------------------
\section{Antecedentes}
\label{sec:antecedentes}

% Base: sección B del anteproyecto.
Las limitaciones de muestreo del espectro angular basado en FFT se
identificaron tempranamente \cite{li2007}. El método de banda limitada
\cite{matsushima2009} restringe el cálculo a las frecuencias que sí cumplen la
condición de muestreo y es hoy la solución más difundida; una alternativa
opuesta amplía la ventana de entrada para comprimir el muestreo en frecuencia
\cite{yu2012}; también se han explorado transformadas de Fourier con muestreo
no uniforme, que amplían el rango útil a costa de una implementación más
compleja \cite{kim2014}.

El MPASM se propuso en 2020 junto con una implementación en Python de acceso
abierto \cite{zhao2020, zhao2020code}. Los resultados publicados lo comparan
contra el ASM convencional, el método de banda limitada y la integral de
Rayleigh--Sommerfeld, cuya evaluación directa sirve de referencia porque no
impone restricciones de muestreo sobre el plano de salida. Toda esa
validación, sin embargo, se hizo en simulación ---un haz gaussiano atravesando
una lente---, sin datos experimentales.

En DLHM, el grupo de Óptica de la Universidad Nacional de Colombia (Sede
Medellín) ha abordado la reconstrucción mediante la implementación no
aproximada de la integral de Rayleigh--Sommerfeld (RSD), que elimina las
restricciones sobre el rango de propagación a costa de un elevado costo
computacional, mitigado en GPU \cite{buitrago2019}; por esa ausencia de
restricciones de muestreo se emplea aquí como referencia de exactitud
numérica. La reconstrucción de rutina en el montaje DLHM del laboratorio, en
cambio, se realiza con un propagador de espectro angular que incorpora la
iluminación de onda esférica del \emph{pinhole}, y es esa la línea base frente
a la cual se contrastan los métodos implementados. Este trabajo se inserta en
esa línea, evaluando una familia de propagadores ---la del producto
matricial--- no explorada antes en DLHM.

% ============================================================
% 3. MATERIALES Y MÉTODOS
% Una sección por fase del anteproyecto. Aquí va lo que se
% construyó y cómo se midió; los números van en el capítulo 4.
% Los comentarios "Repo:" señalan de dónde sale cada parte.
% ============================================================
\chapter{Materiales y métodos}
\label{cap:metodos}

\pendiente{un párrafo sobre la metodología en cascada y la validación en dos
niveles (referencias numéricas y hologramas de laboratorio), con todos los
métodos aplicados al mismo dato.}

% ------------------------------------------------------------
\section{Marco computacional unificado (Fase 1)}
\label{sec:implementacion}
% Repo: CamposT/propagadores.py, propagacion.py, backend.py (CPU/GPU),
%       campos.py; scripts/mpasm_minimo.py; código de referencia
%       en referencia/ y el parche de scripts/parchar_referencias.py.

\pendiente{arquitectura del paquete y la interfaz común de los tres
propagadores; soporte de GPU; cómo se reprodujo primero el código de
referencia \cite{zhao2020code}.}

\subsection{Verificación}
\label{sec:verificacion}
% Repo: CamposT/referencias.py, metricas.py; scripts/exactitud.py;
%       tests/test_rs1.py, test_propagadores.py.

\pendiente{referencia de Rayleigh--Sommerfeld en 1D; haz gaussiano analítico;
métrica SAM \cite{zhao2020}; pruebas automáticas.}

% ------------------------------------------------------------
\section{Calibración numérica (Fase 2)}
\label{sec:calibracion}
% Repo: scripts/comparacion.py, kf_evolucion.py, contraste_referencias.py;
%       resultados/kf, escalado_N.csv, escalado_s.csv.

\pendiente{parámetros barridos (distancias, paso y tamaño del sensor, $\lambda$,
$s$, $K_f$); cómo se construyen los mapas de error; criterio de «precisión
aceptable».}

\subsection{Iluminación esférica y distancia efectiva}
\label{sec:esferica}
% Repo: CamposT/montaje.py, retropropagacion.py; tarea 24 del cronograma.

\pendiente{el cambio de coordenadas antes de propagar: paso $\delta/M$ y la
distancia efectiva $z_\text{eq}$, validados contra RS.}

\subsection{Contraste con la transformada Z \emph{chirp}}
\label{sec:czt}

\pendiente{cómo se implementó la CZT y cómo se comparó con el MPASM.}

% ------------------------------------------------------------
\section{Montaje y registro experimental (Fase 3)}
\label{sec:montaje}
% Repo: docs/hoja_parametros_montaje.md; CamposT/montaje.py.

\pendiente{esquema del montaje; tabla de parámetros medidos ($\lambda$,
$\delta_x$, $\delta_y$, $N_x \times N_y$, $L$, $z$); muestras (USAF 1951 u
otra) y protocolo de adquisición.}

% ------------------------------------------------------------
\section{Reconstrucción y métricas de comparación (Fase 4)}
\label{sec:metricas}
% Repo: CamposT/pipeline.py, roi.py, metricas.py; scripts/barrido_foco.py,
%       retro_*.py, mi_prueba_*.py.

\pendiente{cadena de reconstrucción; búsqueda del plano de foco por nitidez
(sin verdad de terreno); resolución lateral sobre el \emph{target}, contraste,
relación señal--ruido y tiempo de cómputo.}

% ============================================================
% 4. RESULTADOS Y ANÁLISIS
% Organizado por objetivo específico, para que se vea que cada
% uno se cumplió. Cada sección cierra con qué dice el resultado,
% no solo con la figura.
% ============================================================
\chapter{Resultados y análisis}
\label{cap:resultados}

% ------------------------------------------------------------
\section{Verificación de la implementación (objetivo 1)}
\label{sec:res-obj1}

\pendiente{reproducción del código de referencia; error frente a RS y al haz
gaussiano; tiempos en CPU y GPU (resultados/tiempos.png,
resultados/aceleracion.png).}

% ------------------------------------------------------------
\section{Validez bajo el muestreo de DLHM (objetivo 2)}
\label{sec:res-obj2}

\pendiente{mapas de error por método; rangos de distancia, apertura numérica,
$s$ y $K_f$ donde cada uno es confiable; validación de la iluminación
esférica; comparación MPASM--CZT.}

% ------------------------------------------------------------
\section{Reconstrucción de hologramas experimentales (objetivo 3)}
\label{sec:res-obj3}

\pendiente{reconstrucciones con cada propagador frente al ASM con onda
esférica del laboratorio; tabla de métricas; figuras del \emph{target}.}

% ------------------------------------------------------------
\section{Discusión}
\label{sec:discusion}

\pendiente{responder la pregunta de investigación: qué del control de muestreo
mejora la reconstrucción y en qué condiciones de registro conviene cada
método. Incluir los límites de lo que se midió.}

% ============================================================
% 5. CONCLUSIONES Y TRABAJOS FUTUROS
% ============================================================
\chapter{Conclusiones y trabajos futuros}
\label{cap:conclusiones}

\section{Conclusiones}

\pendiente{una conclusión por objetivo específico, con el número que la
sostiene, y una que responda la pregunta de investigación.}

\section{Trabajos futuros}

\pendiente{lo que quedó fuera del alcance o abierto por los resultados.}


% ---------- Bibliografía ----------
\printbibliography[heading=bibintoc, title={Bibliografía}]

\end{document}
